% Exact candidate proof for independent and production-Rethlas audit.
% This file is isolated from the live book.  It makes no novelty or priority
% claim and must not be promoted before a final/correct production receipt.

\begin{proof}
We use the following exact source facts.

First, Mesquita--Piccione identify rational PL functions on \(X^\beth\)
with the direct limit of invariant vertical \(\mathbf Q\)-Cartier divisors
on dominating analytic test configurations
\cite[Proposition~2.4.3 and Theorem~2.6.4]{MP24}.  Their
Definition~4.1.2 declares such a
function \(\alpha\)-psh precisely when the class
\(\alpha_{\mathcal X}+[D]\) is relatively nef, and their
Definition~4.3.1 defines the full cone by pointwise decreasing nets in
that rational PL model cone, excluding the identically
\(-\infty\) limit.

Second, Boucksom--Jonsson define PL positivity by the formally analogous
algebraic relative-nef condition on algebraic test configurations
\cite[Definition~3.2]{BJ22}.  Their Lemma~4.3 identifies this PL
condition with membership in their full psh cone.  Their Theorem~4.5
gives closure under decreasing nets whose limit is generically finite.
Finally, if \(\alpha\) is ample, their Theorem~4.15(iii) writes every
\(\alpha\)-psh function as the pointwise decreasing limit of a net in
\(\mathcal H^{\mathrm{dom}}(\alpha)\).  By Definition~3.8 in that source,
each member of \(\mathcal H^{\mathrm{dom}}(\alpha)\) is a rational PL
function represented on a dominating algebraic test configuration by a
vertical \(\mathbf Q\)-Cartier divisor \(D\) for which
\(\alpha_{\mathcal X}+D\) is relatively ample.

We first compare the rational PL lattices and their positive cones.
Mesquita--Piccione rational PL functions are the
\(\mathbf Q\)-vector space spanned by flag-ideal functions.  Choose the
finitely many flag ideals occurring in one presentation and clear all
denominators.  Each flag ideal on \(X\times\mathbf P^1\) has an expansion

\[
                    \mathfrak a=\sum_j\mathfrak a_jt^j
\]

with only finitely many nontrivial coherent coefficient ideals on \(X\).
Since \(X\) is projective, GAGA algebraizes each \(\mathfrak a_j\), and
the displayed finite weighted expansion algebraizes the flag ideal
together with its \(\mathbf C^*\)-action.  Take one projective model
principalizing all the resulting algebraic flag ideals.  Blowup is
relative Proj of the Rees algebra, so analytification commutes with this
blowup; proper GAGA for finite algebras likewise identifies algebraic and
analytic normalizations.  A \(\mathbf C^*\)-equivariant projective
resolution in characteristic zero then gives a common smooth projective
algebraic test configuration carrying the pulled-back vertical divisors
and their rational linear combination.  Conversely, analytifying an
algebraic flag ideal gives the corresponding analytic flag ideal.  The
restriction from \(\mathbf P^1\) to \(\mathbf A^1\) is a
Boucksom--Jonsson test configuration, while the equivariantly trivial
compactification over \(\mathbf P^1\) reverses this passage; all the
vertical divisors under discussion are supported over \(0\).
The
coefficient formula on vertical components in

\[
                    D=\sum_E b_E\varphi(r(v_E))E
\]

is the same in both theories.  Thus \(\iota\) identifies their rational
PL lattices, and every rational PL function may be tested on one common
smooth projective algebraic model.  Notice that this argument does not
claim that every analytic test configuration is algebraizable.  The
underlying space and ideal evaluations agree by
\cite[Remark~1.3.3]{MP24}.

Let \(\mathcal Y\to X\times\mathbf P^1\) be such a common smooth
projective model, let \(D\) represent \(\varphi\), and put

\[
                         \beta=\alpha_{\mathcal Y}+[D].
\]

We make explicit why MP positivity may be checked on this algebraic
determination.  Suppose an analytic determination
\((\mathcal X,D_{\mathcal X})\) originally represents the same
\(\varphi\).  Cofinality and directedness of projective snc analytic test
configurations \cite[Remark~2.5.2]{MP24} give a common analytic
domination

\[
 \mathcal Z\xrightarrow{p}\mathcal X,
 \qquad
 \mathcal Z\xrightarrow{q}\mathcal Y .
\]

After one further refinement if necessary, injectivity in the direct-limit
isomorphism \cite[Theorem~2.6.4]{MP24} gives

\[
                  p^*D_{\mathcal X}=q^*D.
\]

Thus the two pulled-back cohomological classes on \(\mathcal Z\) agree.
Relative nefness ascends under pullback and descends under these proper
modifications, which is exactly the model-independence assertion in
\cite[Definition~4.1.2 and the proof of Theorem~4.1.5]{MP24}.
Consequently the MP class on \(\mathcal X\) is relatively nef if and only
if the analytified class \(\beta\) on \(\mathcal Y\) is relatively nef.

On a smooth projective variety, the analytic and algebraic nef cones agree
on \(N^1_{\mathbf R}\).  Indeed, algebraic ample classes are Kähler by
Kodaira.  Conversely, let \(\gamma\in N^1_{\mathbf R}\) be Kähler and fix
an ample class \(H\).  Openness of the Kähler cone gives an
\(\varepsilon>0\) such that
\(\gamma-\varepsilon H\) is still Kähler.  Hence
\(\gamma\cdot C\geq\varepsilon H\cdot C\) for every irreducible curve
\(C\).  By positive homogeneity and continuity this inequality holds on
the closed cone of curves, and Kleiman's criterion places \(\gamma\) in
the algebraic ample cone.  Thus the Kähler and ample cones agree inside
\(N^1_{\mathbf R}\), and taking their closures identifies analytic and
algebraic nefness.

Apply this observation to every irreducible component of the
simple-normal-crossings central fibre of \(\mathcal Y\).  Every such
component is smooth and projective.  Restricting the compactified model
to \(\mathbf A^1\) gives the BJ test configuration, while the canonical
equivariantly trivial compactification recovers the MP model.  On every
noncentral fibre, including the fibre over infinity in the compactified
model, \(\beta\) restricts to the ample class \(\alpha\).  Therefore
\(\beta\)
is relatively nef in the analytic sense used by Mesquita--Piccione if and
only if it is relatively nef in the algebraic sense used by
Boucksom--Jonsson; compare the algebraic central-fibre criterion
\cite[Lemma~3.1]{BJ22}.  The common-refinement paragraph applies this
criterion even when the original analytic representative of \(\varphi\)
was carried by a nonalgebraizable test configuration.  This proves
\eqref{eq:p001-mp-bj-model-cone-candidate}.

We now pass to decreasing completions.  Suppose first that
\(\varphi\in\operatorname{PSH}_{\mathrm{MP}}(\alpha)\).  Choose a net
\(\varphi_i\downarrow\varphi\) in the MP rational PL model cone.  By
\eqref{eq:p001-mp-bj-model-cone-candidate} and
\cite[Lemma~4.3]{BJ22}, every \(\varphi_i\) is BJ
\(\alpha\)-psh.  The limit is upper semicontinuous, being a decreasing
limit of continuous functions.  It is also generically finite.  Indeed,
since \(X\) is connected and smooth, it is irreducible.  The PL maximum
principle \cite[Proposition~3.6(iv)]{BJ22} gives

\[
                  \varphi_i(v)\leq\varphi_i(v_{\mathrm{triv}})
                  \qquad(v\in X^{\mathrm{an,triv}}).
\]

After taking the limit, if
\(\varphi(v_{\mathrm{triv}})=-\infty\), then
\(\varphi\equiv-\infty\), contrary to the MP definition.  Thus the limit
is finite at the trivial valuation and hence generically finite by
\cite[Lemma~1.4(ii)]{BJ22}.  The
decreasing-closure theorem
\cite[Theorem~4.5]{BJ22} now yields
\(\varphi\in\operatorname{PSH}_{\mathrm{BJ}}(\alpha)\).

Conversely, let
\(\psi\in\operatorname{PSH}_{\mathrm{BJ}}(\alpha)\).  By
\cite[Theorem~4.15(iii)]{BJ22}, there is a decreasing net

\[
                 \psi_i\in\mathcal H^{\mathrm{dom}}(\alpha),
                 \qquad \psi_i\downarrow\psi.
\]

Each \(\psi_i\) is rational PL and relatively ample on an algebraic
dominating model.  In particular it lies in the common relatively nef PL
cone of
\eqref{eq:p001-mp-bj-model-cone-candidate}.  Since a BJ psh function is
generically finite, \(\psi\not\equiv-\infty\).  The defining decreasing-
net criterion of
\cite[Definition~4.3.1]{MP24} therefore gives
\(\psi\in\operatorname{PSH}_{\mathrm{MP}}(\alpha)\).  This proves
\eqref{eq:p001-mp-bj-full-cone-candidate} and also shows directly that the
identification preserves and reflects pointwise order.

There is no conflict with
\cite[Remark~4.1.3]{MP24}, which says that the MP relatively Kähler model
locus can be larger than the algebraic Fubini--Study locus.  The proof only
uses equality of the rational relatively nef PL cones and the cofinal BJ
regularization by \(\mathcal H^{\mathrm{dom}}(\alpha)\); it does not claim
equality of those two smaller strictly positive loci.

Finally, when \(\alpha=c_1(L)\) for an ample \(\mathbf Q\)-line bundle
and \(\theta\) is a smooth representative of \(c_1(L)\),
\cite[Theorem~3.14]{DXZ23} identifies the full DXZ cone with the full
Boucksom--Jonsson cone on \(X^{\mathrm{an,triv}}\), without an energy
hypothesis.  Composing that order isomorphism with
\eqref{eq:p001-mp-bj-full-cone-candidate} proves the last assertion.
\end{proof}
